\label{ch:appendix-prompts}

This appendix gives representative excerpts of the prompts referenced from
Chapter~\ref{ch:methods}, one per LLM-driven stage
(Section~\ref{sec:methods-phases}). Prompts are assembled programmatically
from task context (scenario, framework, cluster capacity) and, on
refinement iterations, prior feedback or failure records; the excerpts
below show the fixed scaffolding around those variable parts, with
\texttt{[[double-bracket notes]]} marking where per-iteration content is
substituted in. Long, purely repetitive stretches, in particular the
deployment-specification schema's full field-by-field reference, are
elided with \texttt{[\ldots]} and summarised in prose instead. Every call
is logged in full at run time (\texttt{prompt.log} next to the
corresponding \texttt{response.log} in each phase folder,
Section~\ref{sec:methods-artifacts}), so the exact text sent for any
specific iteration in this thesis's experiments is reproducible from the
results directory even where this appendix only excerpts it.


\section{Decision Prompt}

Sent once per refinement iteration, when routing is not already forced
(Section~\ref{sec:methods-decision}). The pointer block names the
iteration folders the agent's own prior \texttt{<CODE>} and
\texttt{<SPEC>} responses live under in conversation history; the context
block is either the previous iteration's benchmark feedback or its
failure record, never both.

\small
\begin{verbatim}
You are a performance optimization strategist for BaxBench iterative
experiments.

After iteration `iteration-004`, decide what to refine **next**
(`iteration-005`) to improve benchmark **goodput** (sustained rate of
*successful* HTTP responses; failed requests do not count).

**Progress**: Iteration 5 of 10 (refinement); 5 more iteration(s)
remain after this one. Choose the path most likely to lift goodput
within the remaining budget.

You may choose **exactly one** path:

1. **`deployment`** - tune Kubernetes deployment parameters only.
   Levers include backend replicas, concurrency, and CPU/memory;
   database replicas and resources/GUCs; PgBouncer `pooler` and
   `read_pooler`; optional Redis `cache`; and pod `placement`. The
   application source code stays unchanged.
2. **`code`** - improve the **application source code** (performance,
   error handling, DB usage, concurrency). New code must pass
   functional tests. The deployment spec stays unchanged in this
   iteration.

## Context
- Scenario: Recipes
- Environment: Go-net-http
- Iteration: iteration-005

[[pointer block: which iteration folder holds the current <CODE> /
<SPEC>, and whether that code currently passes functional tests]]

[[context block: benchmark feedback from iteration-004, OR its
structured failure record, never both]]

## Output format

Return exactly:

<DECISION>
deployment|code
</DECISION>

<RATIONALE>
One short paragraph explaining why this path should help next.
</RATIONALE>

Use `deployment` or `code` (lowercase) inside `<DECISION>`.
\end{verbatim}
\normalsize


\section{Code Generation}
\label{sec:appendix-prompts-code}

\subsection{Baseline}

Iteration~0 codegen wraps BaxBench's existing scenario prompt (API
contract, target framework, expected output format) in an experiment
preamble stating the goal, iteration progress, and the goodput
optimisation objective:

\small
\begin{verbatim}
You are an application engineer for BaxBench iterative Kubernetes
performance experiments.

## Goal
Implement the application for iteration `iteration-000` so it
**passes functional tests** and performs well once deployed.

**Progress**: Iteration 0 of 10 (baseline)

**Optimization objective**:
Maximize **goodput** (successful HTTP responses per second). The
app runs in a container on Kubernetes under adaptive Locust load
with many concurrent users. Failed requests do not count toward
your score. Write for sustained throughput - efficient DB access,
safe schema init at startup, connection pooling - not just minimal
correctness.

[[one-line scenario performance guidance, e.g. "expect thousands of
concurrent users"]]

## Context
- Scenario: Recipes
- Environment: Go-net-http (listen port 8080)
- Database required: True
- Iteration: iteration-000

[[BaxBench scenario prompt: OpenAPI spec, framework template,
<FILEPATH>/<CODE> output format]]
\end{verbatim}
\normalsize

If a prior attempt within this same iteration failed its functional
tests, the same prompt is resent with a ``your previous attempt failed,
fix it'' block appended (Section~\ref{sec:methods-code}), naming the
specific failing tests and the evidence collected for each.

\subsection{Refinement}

Code refinement replaces the goal paragraph and adds an artifact
pointer instead of inlining the previous source tree:

\small
\begin{verbatim}
## Goal
Improve the **application source code** for iteration
`iteration-006` using the starting codebase referenced below. New
code must pass functional tests. The **deployment spec stays
unchanged** in this iteration.

**Progress**: Iteration 6 of 10 (refinement); 4 more iteration(s)
remain after this one. Budget your remaining iterations accordingly
- pick the changes most likely to lift goodput within what is left.

Keep the same API contract and scenario requirements. Output a
complete replacement codebase using the same `<FILEPATH>` / `<CODE>`
format as initial generation. Use the referenced `<CODE>` response
as your starting point.

- **Application code**: see your `<CODE>` response from
  `iteration-004` in conversation history (passed functional tests).

[[BaxBench scenario prompt, same as baseline]]
\end{verbatim}
\normalsize

A pathology the previous iteration's feedback flagged (for example, idle
read replicas because the code never queries \texttt{DB\_READ\_HOST})
reaches this prompt only indirectly, through the benchmark feedback the
decision stage already saw and the agent's own conversation history, not
as a separately re-injected block.


\section{Deployment Specification}
\label{sec:appendix-prompts-spec}

The spec prompt is the longest of the three: it documents every schema
field's semantics and every hard scheduling rule
(Table~\ref{tab:spec-schema}), but, deliberately, no recommended numeric
value.

\small
\begin{verbatim}
You are a Kubernetes deployment tuning expert for BaxBench
performance experiments.

## Goal
You are refining deployment parameters for iteration `iteration-007`.

**Progress**: Iteration 7 of 10 (refinement); 3 more iteration(s)
remain after this one. Plan your remaining budget - bold experiments
early, consolidate refinements toward the end.

**Optimization objective**: Maximize **goodput** (sustained rate of
*successful* HTTP responses). Failed requests do not count. Raw
throughput with high error rates is NOT a win.

[[one-line scenario performance guidance]]

## Context
- Scenario: Recipes
- Environment: Go-net-http (listen port 8080)
- Iteration: iteration-007

## Conversation history
- **Kubernetes deployment**: see your `<SPEC>` response from
  `iteration-005` in conversation history.

[[validation feedback block: only present on a retry after a parse
or validation error, quoting the exact error(s)]]

## Cluster capacity
Schedulable **workers only** (control-plane excluded). Use
**requests** for scheduling fit.

**Cluster budget (sum across workers, after 15% reserve):**
- CPU: 14280m (~14.3 cores)
- Memory: ~54.4 Gi

**Per-worker capacity (each pod must fit on ONE of these nodes):**
- `node1`: allocatable 4000m CPU / 16.0 Gi memory (schedulable;
  leave ~10% headroom per node)
[[... one line per worker ...]]

[[cluster capacity restated as JSON, a machine-readable mirror of
the bullets above]]

## Scheduling rules (critical - hard limits enforced before deploy)
1. **One pod, one node**: each pod's **requests** must fit entirely
   on at least one worker.
2. **Cluster budget**: sum of all pod requests (backends + primary +
   read replicas + poolers + caches) must fit cluster capacity after
   reserve.
[[... rules 3-5: placement, DB connection budget, framework-managed
fields the agent must not set ...]]

## Spec fields (semantics - you choose values)
The framework validates feasibility; it does not prescribe tuning
targets.

[[... full field-by-field reference for backend, pooler,
read_pooler, cache, database, and database.tuning; condensed in
the deployment-specification schema table (Table 3.4) ...]]

## Tuning heuristics (from benchmark feedback)
- **Primary saturated, replicas idle** -> app must use
  `DB_READ_HOST` for reads before adding `read_pooler` or more
  replicas.
- **PgBouncer `cl_waiting` > 0** in diagnostics -> raise
  `default_pool_size`, `max_client_conn`, or pooler
  `replicas`/`resources`.
[[... remaining heuristics: overload vs. latency reading, low-CPU
bottleneck, one-structural-change-per-iteration guidance ...]]

## Output format
Return exactly one block:

<SPEC>
backend:
  replicas: <int>
  resources:
    cpu_request: <quantity>
    cpu_limit: <quantity>
    memory_request: <quantity>
    memory_limit: <quantity>
[[... full YAML skeleton: optional pooler, read_pooler, cache,
database (incl. tuning, placement); condensed in
the deployment-specification schema table (Table 3.4) ...]]
</SPEC>
\end{verbatim}
\normalsize

The baseline prompt is structurally identical, with the goal paragraph
replaced by ``propose deployment parameters \ldots{} so the application
can sustain very high concurrent user load'' and no conversation-history
pointer block, since there is no prior spec yet to point at.
